\section{Proofs}\label{app: proofs}

\subsection{Derivation of the investor's modified problem}\label{app: proof_investor_problem}

We start by providing a derivation of the investor's modified problem. 
\begin{proof}
    First, we adopt a change of variables and write the investor's problem as follows
    \begin{equation}\label{eq: app_investor_problem}
        V_{i,t} = \max_{\{\tilde{C}_{i,t}, h_{i,t}, B_{i,t}, S_{i,t}\}} (1-\beta) U\left(\tilde{C}_{i,t} \right) + \beta U\left( \Psi^ {-1}\left( \mathbb{E}_{i,t}\left[ \Psi\left( U^ {-1}(V_{i,t+1}\right)\right] \right)  \right),
    \end{equation}
    subject to 
    \begin{equation}\label{eq: app_flow}
            \tilde{C}_{i,t} + Q_t S_{i,t} + B_{i,t} = R_{e,t} Q_{t-1} S_{i,t-1} + R_{b,t} B_{i,t-1} + W_t h_{i,t} -\xi_t \frac{h_{i,t}^{1+\nu}}{1+\nu},
    \end{equation}
    and the natural borrowing
    \begin{equation}\label{eq: app_natural_limit}
        (Q_t+\pi_t) S_{i,t-1} + R_{b,t} B_{i,t-1} + W_t h_{i,t} -\xi_t \frac{h_{i,t}^{1+\nu}}{1+\nu} \geq -\mathcal{H}_{i,t}
    \end{equation}
    It is immediate that the optimal value of $h_{i,t}$ satisfies
    \begin{equation}\label{eq: app_optimal_h}
        W_t = \xi_t h_{i,t}^\nu.
    \end{equation}
    
    We show next that, given $h_{i,t}$ satisfying \eqref{eq: app_optimal_h}, if the sequence $(\tilde{C}_{i,t},B_{i,t}, S_{i,t})$ satisfies \eqref{eq: app_flow} and \eqref{eq: app_natural_limit}, then     there exists $(N_{i,t}, \omega_{i,t})$ such that $(\tilde{C}_{i,t}, N_{i,t}, \omega_{i,t}) $ satisfies \eqref{eq:lom_N} and $N_{i,t} \geq 0$. 
    Conversely, if $(\tilde{C}_{i,t}, N_{i,t}, \omega_{i,t}) $ satisfies \eqref{eq:lom_N} and $N_{i,t} \geq 0$, there exists $(B_{i,t}, S_{i,t})$ such that  $(\tilde{C}_{i,t},B_{i,t}, S_{i,t})$         satisfies \eqref{eq: app_flow} and \eqref{eq: app_natural_limit}.
    
    From the definition of the return on human wealth, we have that $W_t h_{i,t} -\xi_t \frac{h_{i,t}^{1+\nu}}{1+\nu} = R_{h,t-1} \mathcal{H}_{i,t-1} - \mathcal{H}_{i,t}$, which allow us to write \eqref{eq: app_flow} and \eqref{eq: app_natural_limit} as follows:
    \begin{equation}\label{eq: app_flow_N}
        \tilde{C}_{i,t} + Q_t S_{i,t} + B_{i,t} + \mathcal{H}_{i,t} = N_{i,t}, \qquad \qquad N_{i,t} \geq 0.
    \end{equation}
    
    % As hours are equalized across investors, then the return on human wealth is same for all investors, allowing us to write $R_{h_i,t} = R_{h,t}$.
    We consider next the law of motion of total wealth:\small
    \begin{align}
        N_{i,t+1} =\left[ R_{e,t+1} \frac{Q_{t} S_{i,t}}{Q_{t} S_{i,t}+B_{i,t}+\mathcal{H}_{i,t}} + R_{b,t+1} \frac{B_{i,t}}{Q_{t} S_{i,t}+B_{i,t}+\mathcal{H}_{i,t}}+ R_{h_i,t+1} \frac{\mathcal{H}_{i,t}}{Q_{t} S_{i,t}+B_{i,t}+\mathcal{H}_{i,t}}\right] \left( N_{i,t} -\tilde{C}_{i,t}\right).
    \end{align}\normalsize
    
    As markets are dynamically complete, there exists replicating portfolios $(\omega_{h_i,t}, \omega_{e,t})$ such that 
    \begin{equation}
        R_{k,t+1} = \omega_{k,t} R_{r,t+1} + (1-\omega_{k,t}) R_{b,t+1},
    \end{equation}
    for $k\in\{h_i,e\}$.
    
    Combining the previous two conditions, we obtain\small
    \begin{align}
        N_{i,t+1} =\left[ R_{r,t+1} \frac{\omega_{e,t} Q_{t} S_{i,t}+\omega_{h_i,t} \mathcal{H}_{i,t}}{N_{i,t}-\tilde{C}_{i,t}} + R_{b,t+1} \frac{B_{i,t}+(1-\omega_{e,t}) Q_t S_{i,t}+(1-\omega_{h_i,t}) \mathcal{H}_{i,t}}{N_{i,t}-\tilde{C}_{i,t}}\right] \left( N_{i,t} -\tilde{C}_{i,t}\right).
    \end{align}\normalsize
    
    Using the first condition in \eqref{eq: app_flow_N} to solve for $B_{i,t}$, we obtain
    \begin{equation}
        N_{i,t+1} =\left[ \left(R_{r,t+1} - R_{b,t+1}\right) \omega_{i,t} + R_{b,t+1} \right] \left( N_{i,t} -\tilde{C}_{i,t}\right),
    \end{equation}
    where $\omega_{i,t} \equiv \frac{\omega_{e,t} Q_{t} S_{i,t}+\omega_{h_i,t} \mathcal{H}_{i,t}}{N_{i,t}-\tilde{C}_{i,t}}$.
    
    
    
    \end{proof}

\subsection{Proof of Lemma~\ref{lemma: policy_functions}} \label{app: proof_policy_functions}

The next lemma characterizes the value function, the consumption function, and the Euler equations of each investor.

\begin{lemma}[Consumption and Euler equations]\label{lemma: policy_functions}
The household's value function takes the form: 
$V_{i}(N,X,s)=U\left(v_{i}(X,s)N\right),$
where $v_{i}(X,s)$ denotes the \textit{wealth multiplier}. The consumption-wealth
ratio $c_{i}(X,s)=\frac{\tilde{C}_{i}(N,X,s)}{N}$ and Euler equations
for investor $i\in\mathcal{I}$ are given by
\begin{enumerate}[(i)]
\item Consumption-wealth ratio:
\begin{equation}
c_{i}(X,s)=\frac{(\beta^{-1}-1)^{\psi}\mathcal{R}_{i}(X,s)^{1-\psi}}{1+(\beta^{-1}-1)^{\psi}\mathcal{R}_{i}(X,s)^{1-\psi}},\label{eq: c_n}
\end{equation}
where $\mathcal{R}_{i}(X,s)\equiv\Psi^{-1}\left(\mathbb{E}_{i}\left[\Psi\left(v(X',s')R_{i}(X,s,s')\right)|X\right]\right)$.
\item Euler equation for an asset $j\in\{r,b\}$: 
\begin{equation}
1=\mathbb{E}_{i}\left[\left.\Lambda_{i}(X,s,s')R_{j}(X,s,s')\right.\right],\label{eq: eulers_app}
\end{equation}
where, for $\theta\equiv\frac{1-\gamma}{1-\psi^{-1}}$, the investor's SDF is given by 
\begin{equation}
\Lambda_{i}(X,s,s')=\beta^{\theta}\left(\frac{c_{i}(\chi(X,s,s'),s')N'}{c_{i}(X,s)N}\right)^{-\frac{\theta}{\psi}}R_{i}(X,s,s')^{-(1-\theta)}.
\end{equation}
\item The wealth multipliers satisfy:
\begin{equation}
v_{i}(X,s)=U^{-1}\left[U\left(c_{i}(X,s)\right)+\beta U\left(\mathcal{R}_{i}(X,s)\left(1-c_{i}(X,s)\right)\right)\right].\label{eq:value.recursion}
\end{equation}
\end{enumerate}
\end{lemma}
% The value function here equals the static utility evaluated at $v_{i}(X,s)N$.\footnote{
%     The wealth multiplier $v_{i}(X,s)$ is a measure of welfare: $v_{i}(X,s)N$ corresponds to the constant consumption level that achieves the investor's expected utility. 
%     This term is analogous to \cite{lucas1987models} measure of the welfare cost of business cycles.}
%     % ---the constant level of consumption that achieves the same utility as in the original stochastic economy.} 
% The consumption-wealth ratio in (\ref{eq: c_n})
% coincides with the consumption-wealth ratio of a deterministic problem
% with return $\mathcal{R}_{i}(X,s)$. In turn, $\mathcal{R}_{i}(X,s)$
% is a risk-adjusted portfolio return with states weighted by a multiplier,
% $v(X',s')$, that solves the recursion in Equation (\ref{eq:value.recursion}). The effect of $\mathcal{R}_{i}$ on consumption depends
% only on the EIS coefficient, $\psi$, that captures substitution and
% income effects. When $\psi=1$, income and substitution effects cancel
% out; the consumption-wealth ratio is $c_{i}(X,s)=1-\beta$. The following
% section derives several results under this parameterization.


\begin{proof}
First, we verify that the value function takes the form \eqref{eq: value_function}. Given the conjecture about the value function, the Bellman equation for investor $i$ can be written as
\begin{equation}
    \frac{(v_i(X,s) N)^{1-\psi^{-1}}-1}{1-\psi^{-1}} = \max_{\tilde{c}_{i},\omega_{i}}  (1-\beta) \frac{(\tilde{c}_i N)^{1-\psi^{-1}}-1}{1-\psi^{-1}} + \beta \frac{\mathbb{E}_i\left[ (v_i(X',s')N')^{1-\gamma}\right]^{\frac{1-\psi^{-1}}{1-\gamma}}-1}{1-\psi^{-1}},
\end{equation}
subject to $N' = R_{i}(X,s,s') (1-\tilde{c}_i) N$ and $N'\geq 0$.

The first-order conditions for the consumption-wealth ratio and the portfolio share are given by 
\begin{align}
    (1-\beta) \tilde{c}_i^{-\psi^{-1}} &= \beta \mathcal{R}_i(X,s)^{1-\psi^{-1}} (1-\tilde{c}_i)^{-\psi^{-1}}\\
    0& = \mathbb{E}_i\left[ (v_i(X',s') R_{i,n}(X,s,s'))^{-\gamma} v_i(X') (R_{r}(X,s,s') - R_b(X,s))  \right]
\end{align}
where $\mathcal{R}_i(X,s) = \mathbb{E}_i\left[ (v_i(X',s') R_{i}(X,s,s'))^{1-\gamma} |X,s\right]^{\frac{1}{1-\gamma}}$.

Given $\mathcal{R}_i(X,s)$, we can solve for the consumption-wealth ratio:
\begin{equation}
    \tilde{c}_i(X,s) = \frac{ (\beta^{-1}-1)^{\psi} \mathcal{R}_i(X,s)^{1-\psi}}{1+ (\beta^{-1}-1)^{\psi} \mathcal{R}_i(X,s)^{1-\psi}}.
\end{equation}

The envelope condition with respect to $N$ is given by
\begin{equation}\label{eq: app_envelope}
    v_i(X)^{1-1/\psi} = \beta \mathcal{R}_i(X) ^{1-1/\psi}  (1-\tilde{c}_i(X))^{-1/\psi} \Rightarrow 
    \tilde{c}_i(X) = (1-\beta)^{\psi} v_i(X)^{1-\psi}.
\end{equation}

From the optimality condition for the risky asset, we obtain 
\begin{equation}\label{eq: app_optimality_omega}
    \mathbb{E}_i\left[  (v_i(X',s') R_{i}(X,s,s'))^{1-\gamma}  \right] = \mathbb{E}_i\left[  v_i(X',s')^{1-\gamma} R_{i}(X,s,s')^{-\gamma} R_j(X,s,s') \right],
\end{equation}
for $j \in \{r,b\}$.

Raising the envelope condition \eqref{eq: app_envelope} to the power $\theta \equiv \frac{1-\gamma}{1-\psi^{-1}}$, using the definition of $\mathcal{R}_i(X)$ and condition \eqref{eq: app_optimality_omega}, we obtain
\begin{equation}
    1 = \mathbb{E}_i\left[ \beta^\theta \left(\frac{v_i(X',s')}{v_i(X,s')}\right)^{1-\gamma} R_{i}(X,s,s')^{-\gamma} R_j(X,s,s') (1-\tilde{c}_i(X,s))^{-\theta/\psi } \right].
\end{equation}

Using the condition $v_i(X) = (1-\beta)^{\frac{1}{1-\psi^{-1}}} \tilde{c}_i(X)^{-\frac{\psi^{-1}}{1-\psi^{-1}}}$, we obtain the Euler equations
\begin{equation}
    1 = \mathbb{E}_i\left[ \beta^\theta \left(\frac{\tilde{c}_i(X',s') N'}{\tilde{c}_i(X,s) N}\right)^{-\frac{\theta}{\psi}} R_{i}(X,s,s')^{-(1-\theta)} R_j(X,s,s') \right].
\end{equation}

This concludes the derivation of the consumption-wealth ratio and the Euler equations for the two assets. It remains to check that the value function takes the form \eqref{eq: value_function}, which amounts to show that $v_i(X)$ indeed does not depend on $N$. Notice that $\tilde{c}_i(X,s)$ and $\omega_i(X,s)$ do not depend on $N$. We can then write the Bellman equation as follows:
\begin{equation}
    v_i(X,s)^{1-\psi^{-1}} =  (1-\beta) \tilde{c}_i(X,s)^{1-\psi^{-1}} + \beta \mathbb{E}_i\left[ (v_i(X',s')R_{i}(X,s,s')(1-\tilde{c}_i(X)))^{1-\gamma}\right]^{\frac{1-\psi^{-1}}{1-\gamma}},
\end{equation}
for $\psi \neq 1$ and
\begin{equation}
    \log v_i(X,s) = (1-\beta) \log \tilde{c}_i(X,s) + \beta \log \mathbb{E}_i\left[ (v_i(X',s')R_{i}(X,s,s')(1-\tilde{c}_i(X,s)))^{1-\gamma}\right]^{\frac{1}{1-\gamma}}.
\end{equation}
which is independent of $N$, which confirms our conjecture for the value function \eqref{eq: value_function}.


\end{proof}


\subsection{Proof of Lemma \ref{lemma: portfolio_share}}\label{app: proof_portfolio_share}
    
The following lemma characterizes households' portfolio
weight in terms of the economy-wide SDF, the market
prices, and their beliefs.

\begin{lemma}[Portfolio share]\label{lemma: portfolio_share} The
shares of total wealth invested in the risky asset are 
\begin{equation}
\omega_{i}(X,s)=\frac{1}{\Delta R_{r}(X,s)}\left[\frac{\tilde{p}_{i}(X,s,H)}{p_{s,H}\Lambda(X,s,H)}-\frac{\tilde{p}_{i}(X,s,L)}{p_{s,L}\Lambda(X,s,L)}\right],\label{eq: portfolio_share}
\end{equation}
where $\tilde{p}_{i}(X,s,s')$ is 
\[
\tilde{p}_{i}(X,s,s')=\frac{(p_{ss'}^{i})^{\frac{1}{\gamma}}\left[v_{i}(\chi(X,s,s'),s')|R_{r}^{e}(X,s,s')|\right]^{\frac{1}{\gamma}-1}}{\sum_{\tilde{s}'\in\{L,H\}}(p_{s\tilde{s}'}^{i})^{\frac{1}{\gamma}}\left[v_{i}(\chi(X,s,\tilde{s}'),\tilde{s}')|R_{r}^{e}(X,s,\tilde{s}')|\right]^{\frac{1}{\gamma}-1}}.
\]
\end{lemma}
Lemma \ref{lemma: portfolio_share}, describes how portfolio shares
depend on distorted probabilities, $\tilde{p}_{i}(X,s,s')$, and
$p_{ss'}\times\Lambda(X,s,s')$. 
The portfolio share $\omega_{i}(X,s)$ in (\ref{eq: portfolio_share}) is increasing in $p_{sH}^{i}$. 
This means that relatively optimistic investors hold more of the risky asset. 

\begin{proof}
Let $X' =\chi(X,s,s')$. 
The optimal portfolio share satisfies the condition
\begin{equation}
    \frac{p_{sL}^i}{p_{sH}^i} \frac{v_i(\chi(X,s,L), L)^{1-\gamma} }{v_i(\chi(X,s,H), H)^{1-\gamma} }\left( \frac{\omega_i(X,s) R_r^e(X,s,L) + 1  }{\omega_i(X,s) R_r^e(X,s,H) + 1  }\right)^{-\gamma}\frac{|R_r^e(X,s,L)|}{R_r^e(X,s,H)} = 1
\end{equation}

Raising both sides to $-\frac{1}{\gamma}$, we obtain
\begin{equation}
    \left(\frac{p_{sL}^i}{p_{sH}^i}\right)^{-\frac{1}{\gamma}} \frac{v_i(\chi(X,s,L), L)^{1-\frac{1}{\gamma}} }{v_i(\chi(X,s,H), H)^{1-\frac{1}{\gamma}} } \frac{\omega_i(X,s) R_r^e(X,s,L) + 1  }{\omega_i(X,s) R_r^e(X,s,H) + 1 }\frac{|R_r^e(X,s,L)|^{-\frac{1}{\gamma}}}{R_r^e(X,s,H)^{-\frac{1}{\gamma}}} = 1
\end{equation}

Rearranging the expression above, we obtain
\begin{equation}\label{app:omega_returns}
\omega_{i}(X,s) =  \frac{\tilde{p}_i(X,s,H)}{|R_r^e(X,s,L)|} -  \frac{\tilde{p}_i(X,s,L)}{R_r^e(X,s,H)}, 
\end{equation}
where
\begin{equation}
    \tilde{p}_{i}(X,s,s') = \frac{(p_{ss'}^i)^{\frac{1}{\gamma}} \left[ v_{i}(\chi(X,s,s'),s') |R_r^e(X,s,s')| \right]^{\frac{1}{\gamma}-1}}{\sum_{s'\in \{L,H\} } (p_{ss'}^i)^{\frac{1}{\gamma}} \left[ v_{i}(\chi(X,s,s'),s') |R_r^e(X,s,s')| \right]^{\frac{1}{\gamma}-1} }.
\end{equation}

The SDF in this economy is given by
\begin{align}
    \Lambda(X,s,s') &= \frac{1}{p_{ss'}} \frac{1}{R_f(X,s)} \frac{|R_r(X,s,-s')-R_f(X,s)|}{\Delta R_{r}(X,s)},
\end{align}
where $\Delta R_r(X,s) = R_r(X,s,H) - R_r(X,s,L)$.

We can then write $\omega_{i}(X,s)$  as follows
\begin{equation}
    \omega_i(X,s) = \frac{1}{\Delta R_r(X,s)} \left[\frac{\tilde{p}_i(X,s,H)}{p_{s,H} \Lambda(X,s,H)}  -\frac{\tilde{p}_i(X,s,L)}{p_{s,L}\Lambda(X,s,L)} \right].
\end{equation}

\paragraph{Log utility.} The formula above simplifies under the case of log utility. 
First, we obtain $\tilde{p}_i(X,s,s') = p_{ss'}^i$. This allows us to write the portfolio share as follows:
\begin{equation}
    \omega_i(X,s) = \frac{p_{s,H}^i}{|R_r^e(X,s,L)|}  -\frac{p_{s,L}^i}{R_r^e(X,s,H)},
\end{equation}
using the expression for $\Lambda(X,s,s')$ and the fact that $R_r^e(X,s,L) < 0 < R_r^e(X,s,H)$ under log utility.

We can rewrite the expression above as follows:
\begin{equation}
    \omega_i(X,s) = \frac{p_{s,H}^i R_r^e(X,s,H) + p_{s,L}^i R_r^e(X,s,L)}{|R_r^e(X,s,L)|R_r^e(X,s,H)} = \frac{\mathbb{E}_{i,s}[R_r^e(X,s,s')]}{|R_r^e(X,s,L)|R_r^e(X,s,H)}.
\end{equation}

The denominator in the expression above can be written as follows:
\begin{align}
    |R_r^e(X,s,L)|R_r^e(X,s,H) = \frac{\mathcal{L}'(X,s) - x_L}{(1-\alpha)\mathcal{L}'(X,s) } \frac{x_H - \mathcal{L}'(X,s)}{(1-\alpha)\mathcal{L}'(X,s) }.
\end{align}

The risk-neutral expectation of productivity growth can be written as follows:
\begin{equation}
    \mathcal{L}'(X,s) = p_{sH}^{RN}(X,s) x_H + p_{sL}^{RN}(X,s) x_L.
\end{equation}
where $p_{ss'}^{RN}(X,s)$ is the risk-neutral probability of state $s'$ given state $s$.

Combining the previous two expressions we obtain:
\begin{equation}
    |R_r^e(X,s,L)|R_r^e(X,s,H)= p_{sH}^{RN}(X,s) p_{sL}^{RN}(X,s) \Delta R_r^e(X,s)^2,
\end{equation}
where $\Delta R_r^e(X,s) = R_r^e(X,s,H) - R_r^e(X,s,L) = \frac{\Delta x}{(1-\alpha) \mathcal{L}'(X,s)} $. 
The expression above is the variance of the excess return on the risky asset under risk-neutral probabilities.

The risk-neutral probability of state $s'$ given state $s$ is given by
\begin{equation}
    p_{ss'}^{RN}(X,s) =\frac{p_{ss'} \Lambda(X,s,s')}{\mathbb{E}_s[\Lambda(X,s,s')]} = \frac{R_r^e(X,s,-s')}{\Delta R_r^e(X,s)}.
\end{equation}

We can then write the risk-neutral probabilities as follows:
\begin{equation}
    p_{sL}^{RN}(X,s) = \frac{R_r^e(X,s,H)}{\Delta R_r^e(X,s)} = \frac{x_H - \mathcal{L}'(X,s)}{\Delta x }, \qquad 
    p_{sH}^{RN}(X,s) = \frac{R_r^e(X,s,L)}{\Delta R_r^e(X,s)} = \frac{\mathcal{L}'(X,s)-x_L}{\Delta x }.
\end{equation}

\paragraph{Diffusion-like approximation.} To better interpret the expression for the portfolio share, it is useful to consider an approximation analogous to the continuous-time limit for diffusion processes. Given $R_{r}(X,s,s')$,  probabilities $p_{ss'}^i$ for household $i$, and a small parameter $\epsilon >0$, we can find $\mu_{i,r}(X,s)$ and $\sigma_{i,r}(X,s)$ that satisfies the conditions\small
\begin{equation}
    R_{r}^e(X,s,H) = \mu_{i,r}(X,s) \epsilon + \sqrt{\frac{p_{sL}}{p_{sH}}} \sigma_{i,r}(X,s) \sqrt{\epsilon}, \qquad 
    R_{r}^e(X,s,L) = \mu_{i,r}(X,s) \epsilon - \sqrt{\frac{p_{sH}}{p_{sL}}} \sigma_{i,r}(X,s) \sqrt{\epsilon}, 
\end{equation}\normalsize
which gives us the expected value and variance for household $i$:
\begin{equation}
    \mathbb{E}_i[R_r^e(X,s,s')|X,s] = \mu_{i,r}(X,s) \epsilon, \qquad \qquad 
    Var_i[R_r^e(X,s,s')|X,s] = \sigma^2_{i,r}(X,s) \epsilon.
\end{equation}
Similarly, we can write $R_b(X,s) = 1 + r_b(X,s) \epsilon$. 

From Equation \eqref{app:omega_returns}, and assuming $\gamma = 1$, we obtain
\begin{align}
    \omega_{i}(X,s) &= R_b(X,s) \frac{p_{s,H}^i R_r^e(X,s,H) + p_{s,L}^i R_r^e(X,s,L)}{|R_r^e(X,s,L)|R_r^e(X,s,H)}\nonumber\\
    &=(1+r_b(X,s) \epsilon) \frac{\mu_{i,r}(X,s) \epsilon}{\left(\sqrt{\frac{p_{sH}}{p_{sL}}} \sigma_{i,r}(X,s) \sqrt{\epsilon}- \mu_{i,r}(X,s) \epsilon\right)\left(\mu_{i,r}(X,s) \epsilon + \sqrt{\frac{p_{sL}}{p_{sH}}} \sigma_{i,r}(X,s) \sqrt{\epsilon} \right)},
\end{align}
where we used the fact that $R_r^e(X,s,L) < 0$ by no-arbitrage. 

In general, $(\mu_{i,r}(X,s),\sigma_{i,r}(X,s))$ and $p_{ss'}^i$ are functions of $\epsilon$. 
Assuming that $\mu_{i,r}(X,s) = \mathcal{O}(1)$, $\sigma_{i,r}(X,s)) = \mathcal{O}(1)$, and $p_{ss'}^i = \mathcal{O}(1)$, we can write the expression $\omega_{i}(X,s)$ as follows:\footnote{These assumptions are analogous to the ones used by e.g. \citet{merton1992continuous} to derive the continuous-time limit with diffusion processes. Allowing for rare events, $p_{ss'}^i = \mathcal{O}(\epsilon)$ for some $s'$, would lead to a jump-diffusion process.  }
\begin{equation}
    \omega_{i}(X,s) =  \frac{\mu_{i,r}(X,s)}{\sigma_{i,r}^2(X,s)} + \mathcal{O}(\epsilon).
\end{equation}

\end{proof}


\subsection{Proof of Proposition \ref{prop:E_sharpe}}\label{app: proof_E_sharpe}
\begin{proof}

First, we compute the Sharpe ratio on the risky asset. 
We will compute expectations using the objective measure, but a similar calculation gives the Sharpe ratio using the investors' subjective beliefs.
The expected excess return is given by
\begin{equation}
    \mathbb{E}\left[ R^e_r(X,s,s') \right] = p_{sL} R^e_r(X,s,L) + p_{sH} R^e_r(X,s,H).
\end{equation}

The variance of excess returns is given by
\begin{equation}
    Var[R^e_r(X,s,s')] = p_{sL} p_{sH} \Delta R^e_r(X,s)^2.
\end{equation}

The Sharpe ratio in the risky asset is then given by
\begin{equation}
    \frac{\mathbb{E}[R^e_r(X,s,s')]}{\sqrt{Var[R^e_r(X,s,s')]}} = \sqrt{\frac{p_{sL}}{p_{sH}}} \frac{R_r^e(X,s,L)}{\Delta R_r^e(X,s)} + \sqrt{\frac{p_{sH}}{p_{sL}}} \frac{R_r^e(X,s,H)}{\Delta R_r^e(X,s)}.
\end{equation}

We can write the expression above in terms of the economy's SDF.
The SDF under the objective measure can be written as
\begin{align}
    \Lambda(X,s,L) &= \frac{\mathbb{E}[\Lambda(X,s,s')]}{p_{sL}}  \frac{R_r^e(X,s,H)}{\Delta R_{r}^e(X,s)}, \qquad 
    \Lambda(X,s,H)  = -\frac{\mathbb{E}[\Lambda(X,s,s')]}{p_{sH}} \frac{R_r^e(X,s,L)}{\Delta R_{r}^e(X,s)}.
\end{align}

Combining the expressions above, we obtain
\begin{equation}
    \frac{\mathbb{E}[R^e_r(X,s,s')]}{\sqrt{Var[R^e_r(X,s,s')]}} = \sqrt{p_{sL} p_{sH}} \frac{\Lambda(X,s,L) - \Lambda(X,s,H)}{\mathbb{E}[\Lambda(X,s,s')]}.
\end{equation}

We consider next how the Sharpe ratio affects the risk-neutral expectation of future productivity growth.
The risk-neutral expectation of productivity is given by
\begin{equation}
    \mathbb{E}^Q\left[ x_{t+1} \right] = p_{sL} \frac{\Lambda(X,s,L)}{\mathbb{E}[\Lambda(X,s,s')]} x_L + p_{sH} \frac{\Lambda(X,s,H)}{\mathbb{E}[\Lambda(X,s,s')]} x_H.
\end{equation}

The difference between the expected value of productivity under the physical measure and the risk-neutral measure is given by
\begin{equation}
    \mathbb{E}[x_{t+1}] - \mathbb{E}^Q[ x_{t+1} ] =  p_{sL} \frac{\mathbb{E}[\Lambda(X,s,s')]-\Lambda(X,s,L)}{\mathbb{E}[\Lambda(X,s,s')]} x_L + p_{sH} \frac{\mathbb{E}[\Lambda(X,s,s')]-\Lambda(X,s,H)}{\mathbb{E}[\Lambda(X,s,s')]} x_H.
\end{equation}

Rearranging the expression above, we obtain
\begin{equation}
    \mathbb{E}[x_{t+1}] - \mathbb{E}^Q[ x_{t+1} ] = p_{sL}p_{sH} \frac{\Lambda(X,s,L)-\Lambda(X,s,H)}{\mathbb{E}[\Lambda(X,s,s')]} \Delta x,
\end{equation}
where $\Delta x = x_H-x_L$.

Using the expression for the Sharpe ratio, we obtain
\begin{equation}
    \mathbb{E}^Q[x_{t+1}] = \mathbb{E}[x_{t+1}] - \sqrt{p_{sL}p_{sH}}     \frac{\mathbb{E}[R^e_r(X,s,s')]}{\sqrt{Var[R^e_r(X,s,s')]}} \Delta x.
\end{equation}
Using the fact that $\mathcal{L}'(X_t, s_t) = \mathbb{E}^Q_t[x_{t+1}]$ and $\sigma_s[x_{t+1}] = \sqrt{p_{sL}p_{sH}} \Delta x$, we obtain the result of Proposition \ref{prop:E_sharpe}.
\end{proof}

\subsection{Proof of Proposition \ref{prop:asset_prices}}\label{app: proof_asset_prices}

\begin{proof} We start by deriving the process for returns.  From the market clearing condition for goods, we obtain
\begin{equation}
    \frac{x_s h(\mathcal{L})^\alpha - \xi \frac{h(\mathcal{L})^{1+\nu}}{1+\nu} }{P(X,s)} = 1-\beta
\end{equation}

The return on the surplus claim is given by
\begin{equation}
    R_{p}(X,s,s') = \frac{x_s P(\chi(X,s,s'),s')}{P(X,s) - \left( x_s h(\mathcal{L})^\alpha - \xi \frac{h(\mathcal{L})^{1+\nu}}{1+\nu}\right)} =
    \frac{x_s}{\beta} \frac{x_{s'} h(\mathcal{L}'(X,s))^\alpha - \xi \frac{h(\mathcal{L}'(X,s))^{1+\nu}}{1+\nu} }{x_{s} h(\mathcal{L})^\alpha - \xi \frac{h(\mathcal{L})^{1+\nu}}{1+\nu}}.
\end{equation}

Using the conditions in \eqref{eq: hours_wages}, we can rewrite the expression as follows
\begin{equation}
    R_{p}(X,s,s') = \frac{x_s}{\beta} \frac{x_{s'} \mathcal{L}'(X,s)^\frac{\alpha}{1+\nu-\alpha} - \frac{\alpha}{1+\nu} \mathcal{L}'(X,s)^{\frac{1+\nu}{1+\nu-\alpha}}}{x_{s} \mathcal{L}^\frac{\alpha}{1+\nu-\alpha} - \frac{\alpha}{1+\nu} \mathcal{L}^{\frac{1+\nu}{1+\nu-\alpha}} }.
\end{equation}

Note that the denominator in the expression above is positive if and only if $\mathcal{L} < \frac{1+\nu}{\alpha} x_s$. 
A sufficient condition is given by $\alpha x_H < x_L$, as shown below
\begin{equation}
    \mathcal{L} \leq x_H < \frac{x_L}{\alpha} < \frac{1+\nu}{\alpha} x_s,
\end{equation}
and, similarly, this condition guarantees that the numerator is also positive.

\paragraph{Interest rate.} The interest rate satisfies the condition $R_b(X,s) = \mathbb{E}\left[ \frac{\Lambda(X,s,s')}{\mathbb{E}[\Lambda(X,s,s')]} R_p(X,s,s') \right]$, so $R_b(X,s)$ is given by 
\begin{equation}
    R_{b}(X,s) = \left(1 - \frac{\alpha}{1+\nu}\right)\frac{x_s}{\beta} \frac{ \mathcal{L}'(X,s)^{\frac{1+\nu}{1+\nu-\alpha}}}{x_{s} \mathcal{L}^\frac{\alpha}{1+\nu-\alpha} - \frac{\alpha}{1+\nu} \mathcal{L}^{\frac{1+\nu}{1+\nu-\alpha}} },
\end{equation}
using the fact that $\mathbb{E}\left[ \frac{\Lambda(X,s,s')}{\mathbb{E}[\Lambda(X,s,s')]} x_{s'} \right] = \mathcal{L}'(X,s)$.

The expression above is increasing in $\mathcal{L}'(X,s)$, decreasing in $x_s$, and it is increasing in $\mathcal{L}$ for $s = L$.

\paragraph{Risk premium.} The risk asset's excess return is given by
\begin{equation}
    \frac{R_p(X,s,s')}{R_b(X,s)} = \frac{1}{1-\frac{\alpha}{1+\nu}} \frac{x_{s'} - \frac{\alpha}{1+\nu} \mathcal{L}'(X,s)}{\mathcal{L}'(X,s)}.
\end{equation}

The conditional risk premium is then given by
\begin{equation}
    \mathbb{E}_s[R_{p}^e(X,s,s')] = 
     \frac{1}{1-\frac{\alpha}{1+\nu}} \frac{\mathbb{E}_s[x_{s'}] -\mathcal{L}'(X,s)}{\mathcal{L}'(X,s)},
\end{equation}
given the definition $R_p^e(X,s,s') \equiv \frac{R_p(X,s,s') - R_b(X,s)}{R_b(X,s)}$.

\end{proof}

\subsection{Proof of Proposition \ref{prop: risk} and Corollary \ref{cor: Eprime}}\label{app: proof_risk}

\begin{proof} 
    We start by deriving the expression for the SDF. From the no-arbitrage conditions, we obtain
\begin{equation}
    \Lambda(X,s,s') = \frac{1}{p_{ss'}} \frac{|R_r^e(X,s,-s')|}{\Delta R_r(X,s)}.
\end{equation}

The excess return on the risky asset and the return difference are given by
\begin{equation}
    R_r^e(X,s,s') = \frac{x_{s'} -\mathcal{L}'(X,s)}{(1-\alpha)\mathcal{L}'(X,s)}, \qquad \qquad
    \Delta R_r(X,s)  = \frac{R_f(X,s) \Delta x}{(1-\alpha) \mathcal{L}'(X,s)}.
\end{equation}

Combining the previous two expressions, we obtain
\begin{equation}
    \Lambda(X,s,s') = \frac{1}{p_{ss'}} \frac{|x_{-s'} - \mathcal{L}'(X,s)|}{R_f(X,s) \Delta x}.
\end{equation}

\paragraph{Portfolio share.} Under log utility, the SDF for investor $i$ is given by
\begin{equation}
    \Lambda_i(X,s,s') = \beta \left( R_{i}(X,s,s') (1- (1-\beta)) \right)^{-1} = R_i(X,s,s')^{-1},
\end{equation}
using the fact that consumption growth equals the growth rate of the household's net worth. 

Using the fact that $R_i(X,s,s') = R_f(X,s) \left[1 + \omega_i(X,s) R_r^e(X,s,s') \right]$, and the expression linking the individual and economy-wide SDFs, we obtain
\begin{equation}
    \frac{1}{R_f(X,s)} \frac{1}{1+\omega_i(X,s) R_r^e(X,s,s')} = \frac{1}{p_{ss'}^i} \frac{|x_{-s'} - \mathcal{L}'(X,s)|}{R_f(X,s) \Delta x}.
\end{equation}

Taking the ratio for $s' = H$ and $s' = L$, we obtain
\begin{equation}
    \frac{1+\omega_i(X,s) R_r^e(X,s,L)}{1+\omega_i(X,s) R_r^e(X,s,H)} = \frac{p_{sL}^i}{p_{sH}^i} \frac{\mathcal{L}'(X,s) - x_L}{x_H-\mathcal{L}'(X,s)}.
\end{equation}

Rearranging the expression above, we obtain
\begin{equation}
    p_{sH}^i(x_H-\mathcal{L}'(X,s))\left[1+\omega_i(X,s) R_r^e(X,s,L)\right] = \left[ 1+\omega_i(X,s) R_r^e(X,s,H)\right]p_{sL}^i(\mathcal{L}'(X,s) - x_L).
\end{equation}
Solving for $\omega_i(X,s)$, we obtain
\begin{equation}
    \omega_i(X,s) = \frac{p_{sH}^i(x_H-\mathcal{L}'(X,s)) +p_{sL}^i(x_L - \mathcal{L}'(X,s))}{        p_{sH}^i(\mathcal{L}'(X,s) - x_H) R_r^e(X,s,L) + 
    p_{sL}^i(\mathcal{L}'(X,s) - x_L) R_r^e(X,s,H)} 
\end{equation}

Using the expression for $R_r^e(X,s,s')$, we obtain
\begin{equation}
    \omega_i(X,s) = \frac{p_{sH}^i \frac{x_H-\mathcal{L}'(X,s)}{(1-\alpha) \mathcal{L}'(X,s)} +p_{sL}^i \frac{x_L - \mathcal{L}'(X,s)}{(1-\alpha) \mathcal{L}'(X,s)} }{
        \frac{x_H-\mathcal{L}'(X,s)}{(1-\alpha) \mathcal{L}'(X,s)} \frac{\mathcal{L}'(X,s)-x_L}{(1-\alpha) \mathcal{L}'(X,s)} 
} = \frac{\mathbb{E}_{i,s}[R_r^e(X,s,s')]}{Var_{RN,s}[R_r^e(X,s,s')]}.
\end{equation}

We can write the portfolio share as follows:
\begin{equation}
    \omega_i(X,s) = (1-\alpha) \mathcal{L}'(X,s) \frac{p_{sH}^i(x_H-\mathcal{L}'(X,s)) +p_{sL}^i(x_L - \mathcal{L}'(X,s))}{(x_H - \mathcal{L'}(X,s)) (\mathcal{L}'(X,s) - x_L)}
\end{equation}




\paragraph{Demand for risk.} The demand for risk in this economy is given by
\begin{equation}
    \sum_{i=1}^I \eta_i \sigma_s[R_{i,n}(X,s,s')] =     \sqrt{p_{sL} p_{sH}} \left[\frac{p_{sH}(X,s)}{p_{sH}\Lambda(X,s,H)}  -\frac{p_{sL}(X,s)}{p_{sL}\Lambda(X,s,L)} \right],
\end{equation}
where $p_{ss'}(X,s) = \sum_{i=1}^I \eta_{i,t} p_{ss'}^i$, using the fact that $\sigma_s[R_r(X,s,s')] = \sqrt{p_{sH} p_{sL}} \Delta R_r(X,s)$ and the results in Lemma~\ref{lemma: portfolio_share}.

Using the expression for the SDF, the demand for risk can be written as\small
\begin{equation}
\sum_{i=1}^I \eta_i \sigma_s[R_{i,n}(X,s,s')] = \sigma_s[x_{s'}]\frac{\frac{1+\nu-\alpha}{1+\nu}\frac{x_s}{\beta} }{x_{s} \mathcal{L}^\frac{\alpha}{1+\nu-\alpha} - \frac{\alpha}{1+\nu} \mathcal{L}^{\frac{1+\nu}{1+\nu-\alpha}} } \left[\frac{p_{sH}(X,s) \mathcal{L}'(X,s)^{\frac{1+\nu}{1+\nu-\alpha}}}{\mathcal{L}'(X,s) - x_L}  -\frac{p_{sL}(X,s)\mathcal{L}'(X,s)^{\frac{1+\nu}{1+\nu-\alpha}}}{x_H-\mathcal{L}'(X,s)} \right],    
\end{equation}\normalsize
given $\sigma_s[x_{s'}] = \sqrt{p_{sL} p_{sH}} \Delta x$.

The first term inside brackets in the expression above is decreasing in $\mathcal{L}'(X,s)$ if and only if the following condition holds
\begin{equation}
    \frac{1+\nu}{1+\nu-\alpha} \mathcal{L}'(X,s)^{\frac{1+\nu}{1+\nu-\alpha}-1} (\mathcal{L}'(X,s)-x_L) - \mathcal{L}'(X,s)^{\frac{1+\nu}{1+\nu-\alpha}} < 0 \iff \mathcal{L}'(X,s) < \frac{x_L}{\alpha},
\end{equation}
which holds, given that $\mathcal{L}'(X,s) \leq x_H < \frac{x_L}{\alpha}$.

Therefore, the demand for risk is decreasing in $\mathcal{L}'(X,s)$. 
As $\mathcal{L}'(X,s)$ is decreasing in the Sharpe ratio of the risky asset, then the demand for risk is increasing in the Sharpe ratio.

\paragraph{Supply of risk.} The volatility of returns is given by
\begin{equation}
    \sigma_s[R_p(X,s,s')] = \frac{x_s}{\beta} \frac{\sigma_s[x_{s'}] \mathcal{L}'(X,s)^\frac{\alpha}{1+\nu-\alpha}}{x_{s} \mathcal{L}^\frac{\alpha}{1+\nu-\alpha} - \frac{\alpha}{1+\nu} \mathcal{L}^{\frac{1+\nu}{1+\nu-\alpha}} },
\end{equation}
which is increasing in $\mathcal{L}'(X,s)$ and decreasing in $x_s$. 

% Note that we can write the coefficient of variation of returns as follows:
% \begin{equation}
%     \frac{\sigma_s[R_p(X,s,s')]}{\mathbb{E}_s[R_p(X,s,s')]} = \frac{\sigma_s[x']}{\mathbb{E}_s[x_{s'}]} \frac{\mathbb{E}_s[x_{s'}] \mathcal{L}'(X,s)^\frac{\alpha}{1+\nu-\alpha} }{\mathbb{E}_s[x_{s'}] \mathcal{L}'(X,s)^\frac{\alpha}{1+\nu-\alpha} - \frac{\alpha}{1+\nu} \mathcal{L}'(X,s)^{\frac{1+\nu}{1+\nu-\alpha}} },
% \end{equation}
% so the presence of labor amplifies the volatility of returns. 

\paragraph{Equilibrium.} Combining supply and demand for risk, we obtain\small
\begin{equation}
 \frac{x_s}{\beta} \frac{\sigma_s[x_{s'}] \mathcal{L}'(X,s)^\frac{\alpha}{1+\nu-\alpha}}{x_{s} \mathcal{L}^\frac{\alpha}{1+\nu-\alpha} - \frac{\alpha}{1+\nu} \mathcal{L}^{\frac{1+\nu}{1+\nu-\alpha}} } =  
 \frac{1+\nu-\alpha}{1+\nu}\frac{x_s}{\beta} \frac{\sigma_s[x_{s'}]\mathcal{L}'(X,s)^{\frac{1+\nu}{1+\nu-\alpha}}}{x_{s} \mathcal{L}^\frac{\alpha}{1+\nu-\alpha} - \frac{\alpha}{1+\nu} \mathcal{L}^{\frac{1+\nu}{1+\nu-\alpha}} } \left[\frac{p_{sH}(X,s) }{\mathcal{L}'(X,s) - x_L}  -\frac{p_{sL}(X,s)}{x_H-\mathcal{L}'(X,s)} \right].
\end{equation}\normalsize

The left-hand side is strictly increasing in $\mathcal{L}'(X,s)$, while the right-hand side is strictly decreasing in $\mathcal{L}'(X,s)$ in the interval $x_L < \mathcal{L}'(X,s) < x_H$. 
The right-hand side converges to $+\infty$ as $\mathcal{L}'(X',s)$ approaches $x_L$ from above, and it converges to $-\infty$ as $\mathcal{L}'(X,s)$ approaches $x_H$ from below. 
Therefore, there exists a unique value of $\mathcal{L}'(X,s)$ solving the equation above in this interval. 
Note that the two curves intersect again for $\mathcal{L}'(X,s) > x_H$, which can be seen by noticing that the right-hand is decreasing in $\mathcal{L}'(X,s)$ for $\mathcal{L}'(X,s) > x_H$ and converges to $+\infty$ as $\mathcal{L}'(X,s)$ approaches $x_H$ from above.
Therefore, the economically relevant solution corresponds to the smallest of the two points of intersection.


Rearranging the expression above, we obtain
\begin{equation}
    1 = \frac{1+\nu-\alpha}{1+\nu} \mathcal{L}'(X,s) \frac{p_{sH}(X,s) (x_H-\mathcal{L}'(X,s))- p_{sL}(X,s)(\mathcal{L}'(X,s) - x_L) }{(\mathcal{L}'(X,s) - x_L) (x_H-\mathcal{L}'(X,s))}.
\end{equation}

We then obtain a quadratic equation for $\mathcal{L}'(X,s)$:\small
\begin{equation}
    \frac{\alpha}{1+\nu} \mathcal{L}'(X,s)^2  - \left[ \left( 1 - \frac{1+\nu-\alpha}{1+v} p_{sH}(X,s) \right) x_H + \left( 1 - \frac{1+\nu-\alpha}{1+v} p_{sL}(X,s) \right) x_L  \right] \mathcal{L}'(X,s) + x_L x_H = 0
\end{equation}\normalsize

The equilibrium value is given by the smallest root of the equation above. 

\end{proof}

\subsection{Proof of Proposition~\ref{prop:real_effects_beliefs}}

\begin{proof}
Consider the case where labor can be chosen conditional on the current productivity level.
In this case, hours and detrended profits are given by
\begin{equation}
    h(X,s) = \left( \frac{\alpha x_s}{\xi} \right)^{\frac{1}{1-\alpha}},\qquad
    \pi(X,s) = (1-\alpha) \left( \frac{\alpha}{\xi} \right)^{\frac{\alpha}{1-\alpha}} x_s^{\frac{1}{1-\alpha}}.
\end{equation}
Notice that hours and profits are independent of beliefs. 

Under log utility, returns for the risky asset and the risk-free asset are given by
\begin{equation}
    R_r(X,s,s') = \frac{x_s}{\beta} \frac{\tilde{x}_{s'}}{\tilde{x}_s}, \qquad \qquad 
    R_f(X,s) =\frac{x_s}{\beta}\frac{\tilde{\mathcal{L}}'(X,s)}{\tilde{x}_{s}},
\end{equation}
where $\tilde{\mathcal{L}}'(X,s) \equiv \mathbb{E}\left[ \tilde{x}_{s'} \right]$ and $\tilde{x}_{s} = x_{s'}^{\frac{1}{1-\alpha}}$.

The excess return on the risky asset is then given by
\begin{equation}
    R_r^e(X,s,s') = \frac{\tilde{x}_{s'} - \tilde{\mathcal{L}}(X,s)}{\tilde{\mathcal{L}}(X,s)}.
\end{equation}

The portfolio share is given by
\begin{equation}
    \omega_i(X,s) = p_{sH}^i\frac{\tilde{\mathcal{L}}'(X,s)}{\tilde{\mathcal{L}}'(X,s) - \tilde{x}_L}  -p_{sL}^i\frac{\tilde{\mathcal{L}}'(X,s)}{x_H-\tilde{\mathcal{L}}'(X,s)}.
\end{equation}

The market clearing for the risky asset implies
\begin{equation}
    \overline{p}_{sH}^m \frac{\tilde{\mathcal{L}}'(X,s)}{\tilde{\mathcal{L}}'(X,s) - \tilde{x}_L} = 1 + \overline{p}_{sL}^m\frac{\tilde{\mathcal{L}}'(X,s)}{x_H-\tilde{\mathcal{L}}'(X,s)}.
\end{equation}
The left-hand side is strictly decreasing in $\tilde{\mathcal{L}}'(X,s)$, it approaches $+\infty$ as $\tilde{\mathcal{L}}'(X,s)$ approaches $x_L$ from above, and it approaches $\frac{\overline{p}_{sH}^m x_H}{x_H - x_L}$ as $\tilde{\mathcal{L}}'(X,s)$ approaches $x_H$ from below.
The right-hand side is strictly increasing in $\tilde{\mathcal{L}}'(X,s)$, it approaches $+\infty$ as $\tilde{\mathcal{L}}'(X,s)$ approaches $x_H$ from below, and it approaches $\frac{\overline{p}_{sL}^m x_L}{x_H - x_L}$ as $\tilde{\mathcal{L}}'(X,s)$ approaches $x_L$ from above.
Hence, there exists a unique value of $\tilde{\mathcal{L}}'(X,s)$ solving the equation above.

Therefore, as market beliefs become more optimistic, the interest rate increases and the risk premium decreases.
However, there is no effect on labor demand or return volatility. 
\end{proof}

\subsection{Proof of Proposition \ref{prop: dynamics_beliefs}}\label{app: proof_dynamics_beliefs}

\begin{proof} The portfolio share of investor $i$ is given by
    \begin{equation}
        \omega_i(X,s) = \frac{p_{s,H}^i}{|R_r^e(X,s,L)|}  -\frac{p_{s,L}^i}{R_r^e(X,s,H)},
    \end{equation}
    The return on the investor's portfolio is given by
    \begin{align}
        R_{i}(X,s,s') &= R_f(X,s) \left[ 1 + \omega_i(X,s) R_r^e(X,s,s') \right] \\
        &= R_f(X,s) \left[1 + (x_{s'} - \mathcal{L}'(X,s)) \left( \frac{p_{sH}^i}{\mathcal{L}'(X,s)-x_L}  -\frac{p_{sL}^i}{x_H-\mathcal{L}'(X,s)} \right) \right].
    \end{align}
    
We can write the expression above as follows 
\begin{equation}
    R_i(X,s,L) = \frac{\Delta x R_f(X,s)}{x_H - \mathcal{L}'(X,s)} p_{sL}^i, \qquad  \qquad 
    R_i(X,s,h) = \frac{\Delta x R_f(X,s)}{\mathcal{L}'(X,s) - x_L} p_{sH}^i.
\end{equation}

\paragraph{Wealth share dynamics.} The share of wealth of investor $i$ is given by
\begin{equation}
    \eta_i'(X,s,s') = \frac{\eta_i R_{i,n}(X,s,s')}{\sum_{j=1}^I\eta_j R_{j,n}(X,s,s')} = \frac{\eta_i p_{ss'}^i}{\sum_{j=1}^I \eta_j p_{ss'}^j} = \eta_i \frac{p^i_{ss'}}{p_{ss'}(X)}.
\end{equation}

\paragraph{The case of arbitrary risk aversion.} Consider the case of unit-EIS, so the consumption-wealth ratio is still given by $c_i(X,s) = 1-\beta$, but arbitrary risk aversion $\gamma > 0$. 
In this case, Propostion~\ref{prop: risk} still holds. 
Moreover, $\tilde{\eta}_i(X,s) = \eta_i$. 

As shown in Appendix~\ref{app: proof_portfolio_share}, the portfolio share of investor $i$ is given by
\begin{equation}
    \omega_{i}(X,s) =  \frac{\tilde{p}_i(X,s,H)}{|R_r^e(X,s,L)|} -  \frac{\tilde{p}_i(X,s,L)}{R_r^e(X,s,H)}, 
    \end{equation}
    where
    \begin{equation}
        \tilde{p}_{i}(X,s,s') = \frac{(p_{ss'}^i)^{\frac{1}{\gamma}} \left[ v_{i}(\chi(X,s,s'),s') |R_r^e(X,s,s')| \right]^{\frac{1}{\gamma}-1}}{\sum_{s'\in \{L,H\} } (p_{ss'}^i)^{\frac{1}{\gamma}} \left[ v_{i}(\chi(X,s,s'),s') |R_r^e(X,s,s')| \right]^{\frac{1}{\gamma}-1} }.
    \end{equation}    

A derivation similar to the one above shows that the wealth share of investor $i$ evolves according to
\begin{equation}
    \eta_i'(X,s,s') = \eta_i \frac{\tilde{p}^i_{ss'}}{\tilde{p}^m_{ss'}(X,s)},
\end{equation}
where $\tilde{p}^m_{ss'}(X,s) \equiv \sum_{j=1}^I \eta_j \tilde{p}^j_{ss'}(X,s)$.
Hence, $\frac{\partial \eta_i'(X,s,s')}{\partial \eta_i}  > 0$.

% \paragraph{Long-run wealth dynamics.} Note that the wealth share is a (bounded) martingale under market beliefs:
% \begin{equation}
%     p_{sH}(X) \eta_{i}'(X,s,H) + p_{sL}(X) \eta_{i}'(X,s,L) = \eta_i (p_{sH}^i + p_{sL}^i) = \eta_i.
% \end{equation}

% Therefore, from the martingale convergence theorem, the wealth share of investor $i$ converges.
% This implies that, for every $\epsilon$, there exits $T$ such that 
% \begin{equation}
%     |\eta_{i,T+1} - \eta_{i,T}| < \epsilon \iff \eta_{i} \left| \frac{p_{ss'}^i - p_{ss'}(X)  }{p_{ss'}(X)} \right| < \epsilon,
% \end{equation}
% almost surely, where the economy is at the state $(X,s)$ at period $T$.

% This implies that either $\eta_i$ converges to zero or $p_{ss'}^i$ converges to $p_{ss}(X)$.
% If $p_{ss'}^i \neq p_{ss'}^j$ for any $i,j\in \mathcal{I}$, $i\neq j$, then the wealth share of a single investor converges to one.  



% By definition, $p_{sH}^i > p_{sH}(X)$ for an optimistic investor in state $(X,s)$, then the wealth share of optimists increase in the good state and decline in the bad state. 
% This implies that market beliefs evolve according to
% \begin{equation}
%     p_{s'H}(X') =\sum_{i=1}^I \eta_i' p_{s'H}^i = \sum_{i=1}^I \eta_i \frac{p^i_{ss'}}{p_{ss'}(X)} p_{s'H}^i,
% \end{equation}
% where
% \begin{equation}
%     p_{HH}(X') =\sum_{i=1}^I \eta_i \frac{p_{sH}^i}{p_{sH}(X)} p_{HH}^i \geq p_{sH}(X), \qquad
%     p_{LH}(X') =\sum_{i=1}^I \eta_i \frac{p_{sL}^i}{p_{sL}(X)} p_{LH}^i \leq p_{sH}(X).
% \end{equation}

% This implies that the relative wealth share for investors $i$ and $j$ is given by
% \begin{equation}
%     \frac{\eta_i'(X,s,s')}{\eta_j'(X,s,s')} = \frac{\eta_i}{\eta_j} \frac{p_{ss'}^i}{p_{ss'}^j}.
% \end{equation}

% Suppose investor $j$ beliefs coincides with the objective measure. Then, the ratio above is a martingale:
% \begin{equation}
%     \mathbb{E}_s\left[ \frac{\eta_i'}{\eta_j'} \right] = p_{sL} \frac{\eta_i}{\eta_j} \frac{p_{sL}^i}{p_{sL}} + p_{sH} \frac{\eta_i}{\eta_j} \frac{p_{sH}^i}{p_{sH}} = \frac{\eta_i}{\eta_j}.
% \end{equation}

% If the wealth of investor $j$ is bounded away from zero, then the above martingale is bounded and, from the martingale convergence theorem, it converges almost surely. 
\end{proof}

\subsection{Proof of Corollary \ref{cor:het_dynamics}}\label{app: proof_cor_het_dynamics}
\begin{proof}
Consider an economy that starts at $s = H$ with wealth distribution $\{\eta_i\}_{i=1}^I$ which switches to the low state after either one period (early transition) or two periods (late transition). 
To simplify notation, we write $p_{ss'}(X)$ to denote $\overline{p}^m_{ss'}(X,s)$.

Market beliefs on the low state in the case of an early transition are given by
\begin{equation}
    p_{LH}(X') = \sum_{i=1}^I \eta_i' p_{LH}^i = \sum_{i=1}^I \eta_i \frac{p_{HL}^i}{p_{HL}(X)} p_{LH}^i,
\end{equation}
and market beliefs on the low state in the case of a late transition are given by
\begin{equation}
    p_{LH}(X'') = \sum_{i=1}^I \eta_i'' p_{LH}^i = \sum_{i=1}^I \eta_i' \frac{p_{HL}^i}{p_{HL}(X')} p_{LH}^i,
\end{equation}
where $\eta_i' = \eta_i \frac{p_{HH}^i}{p_{HH}(X)}$. 


Note that if investor $i$ is optimistic,  $p_{HH}^i > p_{HH}(X)$, then $\eta_i' > \eta_i$ and $p_{HL}^{-i}(X') \leq p_{HL}^{-i}(X) $, where $p_{HL}^{-i}(X) \equiv \frac{1}{1-\eta_i}\sum_{j\neq i} \eta_j p_{HL}^j$. 
Hence, the following inequality holds:
\begin{equation}
  \eta_i' \frac{p_{HL}^i}{p_{HL}(X')} =  \frac{\eta_i' p_{HL}^i}{\eta_i' p_{HL}^i + (1-\eta_i') p_{HL}^{-i}(X') } > 
  \frac{\eta_i p_{HL}^i}{\eta_i p_{HL}^i + (1-\eta_i) p_{HL}^{-i}(X) } =  \eta_i \frac{p_{HL}^i}{p_{HL}(X)}. 
\end{equation}

Therefore, there is more weight on the beliefs of investors who were optimistic in the original state in the case of a late transition. In the case of rank-preserving beliefs, these agents are also optimistic in the low state, so the market is more optimistic under a late transition:
\begin{equation}
    p_{LH}(X'') > p_{LH}(X'). 
\end{equation}

Alternatively, the market is now more pessimistic after a late transition in the case of rank-alternating beliefs:
\begin{equation}
    p_{LH}(X'') < p_{LH}(X'). 
\end{equation}

A similar argument shows that, under rank-preserving beliefs, the market is more pessimistic after a late transition when the economy starts at state $s = L$:
\begin{equation}
    p_{HH}(X'') = \sum_{i=1}^I \eta_i' \frac{p_{LH}^i}{p_{LH}(X')} p_{HH}^i < \sum_{i=1}^I \eta_i \frac{p_{LH}^i}{p_{LH}(X)} p_{HH}^i= p_{HH}(X),
\end{equation}
where $\eta_i' = \eta_i \frac{p_{LL}^i}{p_{LL}(X)}$. Alternatively, the market is more optimistic under a late transition in the case of rank-alternating beliefs.

\end{proof}


\subsection{Proof of Lemma~\ref{lem:trading}}\label{app: proof_lemma_trading}
\begin{proof}    
    Turnover is given by
    \begin{equation}
        \tau(X,s, s') = \frac{1}{2} \sum_{i=1}^I | \omega_i(X', s') \eta_i'(X,s,s')  - \omega_i(X,s) \eta_i |,
    \end{equation}
    where $X' = \chi(X,s,s')$ and $\eta'_i(X,s,s') = \eta_i \frac{p_{ss'}^i}{\overline{p}_{ss'}^m(X,s)} $.
    
    \paragraph{Solving for the portfolio share.} Using the expression for the economy-wide SDF and Equation \eqref{eq: portfolio_share}, we can write the portfolio share as follows
    \begin{equation}\label{app:omega_returns}
    \omega_{i}(X,s) =  \frac{p^i_{sH}}{|R_r^e(X,s,L)|} -  \frac{p^i_{sL}}{R_r^e(X,s,H)}. 
    \end{equation}

    The excess return on the risky asset can be written as follows:\small
    \begin{equation}
        R_{r}^e(X,s,s') =  \frac{x_{s'} -\mathcal{L}'(X,s)}{(1 - \alpha) \mathcal{L}'(X,s) }.
    \end{equation}\normalsize

    Combining the previous expressions, we obtain
    \begin{equation}
    \omega_{i}(X,s) = (1-\alpha) \left[p_{sH}^i \frac{ \mathcal{L}'(X,s)}{\mathcal{L}'(X,s) - x_L}- p_{sL}^i \frac{\mathcal{L}'(X,s)}{x_H - \mathcal{L}'(X,s)} \right],
    \end{equation}
    which is strictly decreasing in $\mathcal{L}'(X,s)$ and $\omega_{i}(X,s)  > 1$ if and only if $p_{sH}^i > \overline{p}^m_{sH}(X,s)$. 


    Turnover is then given by
    \begin{equation}\footnotesize
        \tau(X,s,s') = (1-\alpha) \sum_{i=1}^I \eta_i \left\lvert \left(  \frac{ p_{s'H}^i \mathcal{L}'(X',s')}{\mathcal{L}'(X',s') - x_L}-  \frac{p_{s'L}^i\mathcal{L}'(X',s')}{x_H - \mathcal{L}'(X',s')}\right) \frac{p_{ss'}^i}{\overline{p}_{ss'}^m(X,s)}   - \left(  \frac{p_{sH}^i \mathcal{L}'(X,s)}{\mathcal{L}'(X,s) - x_L}-  \frac{p_{sL}^i\mathcal{L}'(X,s)}{x_H - \mathcal{L}'(X,s)}\right)  \right\rvert 
    \end{equation}\normalsize    

    \paragraph{Perturbation.} It is useful to parameterize the dispersion in beliefs as follows:
\begin{equation}
    p_{ss'}^i = p_{ss'}^* + \epsilon \delta_{ss'}^i,
\end{equation}
where $\delta_{sH}^i + \delta_{sL}^i = 0$. 

Notice that all equilibrium variables now depend on $\epsilon$. For instance, the average probability of the high state can be written as
\begin{equation}
    \overline{p}_{sH}^m(X;\epsilon) = p_{sH}^* + \delta_{sH}(X,s) \epsilon  + \mathcal{O}(\epsilon^2),
\end{equation}
where $\delta_{sH}(X,s) \equiv \sum_{i=1} ^I\eta_i \delta_{sH}^i$. 

The risk-neutral expectation of productivity growth is a function of market beliefs, $\mathcal{L}'(X,s; \epsilon) = f(\overline{p}^m_{sH}(X))$, where $f(p)$ satisfies the condition\small
\begin{equation}
    1 = (1-\alpha) \left[ p \frac{f(p)}{f(p)-x_L} - (1-p) \frac{f(p)}{x_H - f(p)}  \right] \Rightarrow 
    f'(p) = \frac{\frac{f(p)}{f(p)-x_L} + \frac{f(p)}{x_H - f(p)}}{p \frac{x_L}{(f(p) - x_L)^2} + (1-p) \frac{x_H}{(x_H-f(p))^2}}.
\end{equation}\normalsize

Let $\mathcal{L}^*(X,s) \equiv \mathcal{L}'(X,s;0)$ denote the value of $\mathcal{L}'(X,s)$ when $\epsilon = 0$. In this case, we can drop the dependence on $X$ and simply write $\mathcal{L}^*(s)$, as $\mathcal{L}'(X,s)$ would only depend on the state $s$. We can then expand $\mathcal{L}'(X,s;\epsilon)$ in $\epsilon$ to obtain:
\begin{equation}
    \mathcal{L}'(X,s;\epsilon) = \mathcal{L}^*(s) + \tilde{\mathcal{L}}(X,s) \epsilon + \mathcal{O}(\epsilon^2),
\end{equation}
where $\tilde{\mathcal{L}}(X,s) = f'(p_{sH}^*) \sum_{i=1}^I \eta_i \delta_{sH}^i$, where $f'(\cdot) > 0$. 

We can then write the portfolio share of investor $i$ as follows
\begin{equation}
    \omega_{i}(X,s; \epsilon) = 1 +\left[ \theta_{\omega,1}(s) \delta_{sH}^i - \theta_{\omega,2}(s) \delta_{sH}(X) \right] \epsilon + \mathcal{O}(\epsilon^2),
\end{equation}
where $\theta_{\omega,1}(s) > 0$ and $\theta_{\omega,2}(s) > 0$
\begin{align}
    \theta_{\omega,1}(s) &\equiv (1-\alpha)\left(  \frac{\mathcal{L}^*(s)}{\mathcal{L}^*(s) - x_L} + \frac{\mathcal{L}^*(s)}{x_H-\mathcal{L}^*(s)} \right)\\
    \theta_{\omega,2}(s) &\equiv (1-\alpha) \left[\frac{ p_{sH}^* x_L}{(\mathcal{L}^*(s) - x_L)^2}+  \frac{p_{sL}^*x_H}{(x_H - \mathcal{L}^*(s))^2} \right] f'(p_{sH}^*).
\end{align}

Using the expression for $f'(\cdot)$, we obtain that $\theta_{\omega,1} = \theta_{\omega,2}$. We can then write $\omega_{i}(X,s;\epsilon)$ as follows:
\begin{equation}
    \omega_{i}(X,s; \epsilon) = 1 +\theta_{\omega,1}(s) \left[ \delta_{sH}^i - \delta_{sH}(X) \right] \epsilon + \mathcal{O}(\epsilon^2),
\end{equation}

The evolution of wealth is given by 
\begin{equation}
    \eta_i'(X,s,s';\epsilon) = \eta_i + \eta_i \frac{\delta_{s s'}^i - \delta_{ss'}(X)}{p_{ss'}^*} \epsilon + \mathcal{O}(\epsilon^2)
\end{equation}

Let $p_H(X,s,s';\epsilon) = \sum_{i=1}^I \eta_i'(X,s,s';\epsilon) p_{s'H}^i$ denote the market-implied probability of the high state after a transition from state $s$ to state $s'$, then
\begin{equation}
    p_H(X,s,s';\epsilon) = p_{s'H}^* + \delta_{s'H}(X)\epsilon + \mathcal{O}(\epsilon^2),
\end{equation}
using the fact that $\sum_{i=1}^I \eta_i \frac{\delta_{ss'}^i - \delta_{ss'}(X)}{p^*_{ss'}} p_{s'H}^* +  \sum_{i=1}^I \eta_i \delta_{s'H}^i =  \sum_{i=1}^I \eta_i \delta_{s'H}^i  \equiv \delta_{s'H}(X)$.

The portfolio share next period is given by
\begin{equation}
    \omega_i'(X,s,s'; \epsilon) = 1 + \theta_{\omega,1}(s') \left[ \delta_{s'H}^i - \delta_{s'H}(X) \right] \epsilon + \mathcal{O}(\epsilon^2).
\end{equation}

Investor $i$'s net purchases of shares is given by
\begin{align}\small
    \Delta S_i(X,s,s'; \epsilon) &= \eta_i\left[\frac{\delta_{s s'}^i - \delta_{ss'}(X)}{p_{ss'}^*}+\theta_{\omega,1}(s')\left( \delta_{s'H}^i - \delta_{s'H}(X) \right) \right]\epsilon\nonumber\\
    &- \theta_{\omega,1}(s) \eta_i\left[ \delta_{sH}^i - \delta_{sH}(X) \right] \epsilon + \mathcal{O}(\epsilon^2)
\end{align}

For simplicity, suppose that investors believe productivity growth to be iid in the reference economy, that is, $p_{Ls'}^* = p_{Hs'}^* $. We can then write the  
\begin{equation}
    \mu_i \Delta S_i(X,s,s';\epsilon) =\left[\underbrace{\Delta \tilde{\omega}_i(X,s,s') \eta_i}_{\text{change-in-beliefs effect}}+
    \underbrace{\Delta \tilde{\eta}_i(X,s,s')}_{\text{rebalancing effect}}\right] \epsilon + \mathcal{O}(\epsilon^2),
\end{equation}
where, given $\theta_{\omega,1}(L) = \theta_{\omega,1}(H) \equiv \kappa_\omega$,
\begin{align}
    \Delta \omega_i(X,s,s') &\equiv \kappa_\omega\left[\left( \delta_{s'H}^i - \delta_{s'H}(X,s) \right) - \left( \delta_{sH}^i - \delta_{sH}(X) \right) \right]\\
    \Delta \eta_i(X,s,s') &\equiv \eta_i \frac{\delta_{s s'}^i - \delta_{ss'}(X)}{p_{ss'}^*}.
\end{align}

\end{proof}

\subsection{Proof of Proposition~\ref{prop:turnover}}\label{app: proof_prop_turnover}

\begin{proof}
    Turnover is given by
\begin{equation}
    \tau(X,s,s';\epsilon) = \frac{1}{2} \sum_{i=1}^I \eta_i\left\lvert  \frac{\tilde{\delta}_{s s'}^i}{p_{ss'}^*} +\kappa_\omega\left(\tilde{\delta}_{s'H}^i - \tilde{\delta}_{sH}^i  \right)\right\rvert \epsilon + \mathcal{O}(\epsilon^2),
\end{equation}
where $\tilde{\delta}_{ss'}^i = \delta_{ss'}^i - \delta_{ss'}(X)$.

Suppose $s = s'$, then
\begin{align}
    \tau(X,s,s';\epsilon) &= \frac{1}{2} \sum_{i=1}^I \eta_i  \frac{\left\lvert \tilde{\delta}_{s s'}^i(X)\right\rvert}{p_{ss'}^*} \epsilon + \mathcal{O}(\epsilon^2)\\
    &= \frac{1}{2} \left[\sum_{i=1}^I \eta_i \frac{\tilde{\delta}_{ss'}^i(X)}{p_{ss'}^*}\textbf{1}_{\tilde{\delta}_{ss'}^i(X)\geq 0} -  \sum_{i=1}^I \eta_i \frac{ \tilde{\delta}_{ss'}^i(X)}{p_{ss'}^*}\textbf{1}_{\tilde{\delta}_{ss'}^i < 0} \right] \epsilon + \mathcal{O}(\epsilon^2)\\
    &= \frac{1}{2} \left[\eta_B \frac{\tilde{\delta}_{ss'}^B(X)}{p_{ss'}^*} +  \eta_S \frac{ |\tilde{\delta}_{ss'}^S(X)|}{p_{ss'}^*} \right] \epsilon + \mathcal{O}(\epsilon^2).
\end{align}
where 
\begin{align}
    \eta_B &\equiv \sum_{i=1}^I \eta_i \textbf{1}_{\tilde{\delta}_{ss'}^i(X) \geq 0}, \qquad \qquad  \tilde{\delta}_{ss'}^B(X) \equiv \frac{1}{\eta_B} \sum_{i=1}^I \eta_i \tilde{\delta}_{ss'}^i(X) \textbf{1}_{\tilde{\delta}_{ss'}^i(X) \geq 0}, \\
    \eta_S &\equiv \sum_{i=1}^I \eta_i \textbf{1}_{\tilde{\delta}_{ss'}^i(X) < 0}, \qquad \qquad  \tilde{\delta}_{ss'}^S(X) \equiv \frac{1}{\eta_S} \sum_{i=1}^I \eta_i \tilde{\delta}_{ss'}^i(X) \textbf{1}_{\tilde{\delta}_{ss'}^i(X) < 0}.
\end{align}

We can write turnover in this case as follows 
\begin{equation}
    \tau(X,s,s';\epsilon) = \eta_B \eta_S \frac{\delta_{ss'}^B(X) -\delta_{ss'}^S(X)}{p^*_{ss'}}\epsilon + \mathcal{O}(\epsilon^2),
\end{equation}
using the fact that $\eta_B \tilde{\delta}_{ss'}^B(X) + \eta_S \tilde{\delta}_{ss'}^S(X) = 0$ and $\tilde{\delta}_{ss'}^B(X) = \eta_S (\delta_{ss'}^B(X) - \delta_{ss'}^S(X))$.    
% $\delta_{ss'}(X) = \eta_B \delta_{ss'}^B(X) + \eta_S \delta_{ss'}^S(X)$.    

\paragraph{Heterogeneous persistence.} We consider next the special case where investors agree about the unconditional mean of $x$, but they disagree about the persistence of the aggregate productivity growth. 

The stationary distribution of beliefs for investor $i$ is given by
\begin{equation}
    \overline{p}_L^i = \frac{p_{HL}^i}{p_{LH}^i+ p_{HL}^i}.
\end{equation}

We assume that $\overline{p}_L^i$ is equalized across investors, so all investors agree about the unconditional mean of $x_{t}$. Note this implies that the likelihood ratio $p_{LH}^i / p_{HL}^i$ is equalized across investors. The unconditional mean is given by
\begin{equation}
    \overline{x} = \frac{p_{HL}^i}{p_{LH}^i+p_{HL}^i} x_L + \frac{p_{LH}^i}{p_{LH}^i+p_{HL}^i} x_H.
\end{equation}

The expected value of $x_{t+1}$ relative to the mean $\overline{x}$ conditional on $x_t = x_L $ is given by
\begin{align}
    \mathbb{E}_i[x_{t+1} - \overline{x}|x_t = x_L] &= p_{LL}^i (x_L - \overline{x})+ p_{LH}^i (x_H - \overline{x})\\
    &= \left[1 +  p_{LH}^i \frac{x_H - x_L}{x_L - \overline{x}} \right](x_L - \overline{x})\\
    &= \left[1 - (p_{LH}^i +p_{HL}^i) \right](x_L - \overline{x}),
\end{align}
using the fact that $\overline{x} - x_L = \frac{p_{LH}^i}{p_{LH}^i+p_{HL}^i}(x_H-x_L)$

We obtain a similar expression conditioning on $x_t = x_H$ instead:
\begin{align}
    \mathbb{E}_i[x_{t+1} - \overline{x}|x_t = x_H] &= p_{HL}^i (x_L - \overline{x})+ p_{HH}^i (x_H - \overline{x})\\
    &= \left[1 -  p_{HL}^i \frac{x_H - x_L}{x_H - \overline{x}} \right](x_H - \overline{x})\\
    &= \left[1 - (p_{LH}^i +p_{HL}^i) \right](x_H - \overline{x}),
\end{align}
using the fact that $x_H - \overline{x} = \frac{p_{HL}^i}{p_{LH}^i+p_{HL}^i}(x_H-x_L)$.

Let $\hat{x}_{t} = x_t - \overline{x}$, we can then write
\begin{equation}
    \mathbb{E}_i[\hat{x}_{t+1}|\hat{x}_t] = \theta_i \hat{x}_t,
\end{equation}
where $\theta_i \equiv 1-  (p_{LH}^i +p_{HL}^i) = p_{HH}^i - p_{LH}^i$.

Given that investors agree about the unconditional mean of $x$, we are able to pin down beliefs as a function of $\theta_i$:
\begin{equation}
    p_{LH}^i = \overline{p}_H (1- \theta_i), \qquad \qquad 
    p_{HH}^i = \overline{p}_H+\overline{p}_L \theta_i.
\end{equation}

\paragraph{Corollary.} Under the assumption investors agree about the unconditional mean of $x_t$, we have that
\begin{equation}
    p_{LH}^i - p_{LH}(X) = - \overline{p}_H (\theta_i - \theta(X)), \qquad
    p_{HH}^i - p_{HH}(X) = \overline{p}_L (\theta_i - \theta(X)),
\end{equation}
where $\overline{\theta}(X) \equiv \sum_{i=1}^I \eta_i \theta_i$.

Notice that we have that  $\tilde{\delta}_{ss'}^i(X) \epsilon = p_{ss'}^i - p_{ss'}(X)$, which gives us
\begin{equation}
    \tilde{\delta}_{LH}^i(X) \epsilon = - p^*_H (\theta_i - \theta(X)), \qquad \qquad
    \tilde{\delta}_{HH}^i(X) \epsilon = (1-p^*_H) (\theta_i - \theta(X)).
\end{equation}

We can then write turnover in the case $s = L$ and $s' = H$ as follows:
\begin{equation}
    \tau(X,L,H;\epsilon) = \frac{1}{2}\left\lvert  \kappa_\omega-1\right\rvert \sum_{i=1}^I \eta_i |\theta_i - \theta(X)| + \mathcal{O}(\epsilon^2).
\end{equation}

Consider now the case $s = H$ and $s' = L$:
\begin{equation}
    \tau(X,H,L;\epsilon) = \frac{1}{2}\left\lvert \kappa_\omega+ 1 \right\rvert \sum_{i=1}^I \eta_i |\theta_i - \theta(X)| + \mathcal{O}(\epsilon^2),
\end{equation}

Suppose now that $s = s' $, then
\begin{align}
    \tau(X,H,H;\epsilon) &= \frac{1}{2}\left\lvert  \frac{p^*_L}{p_{H}^*} \right\rvert \sum_{i=1}^I \eta_i |\theta_i - \theta(X)| \epsilon + \mathcal{O}(\epsilon^2) \\
    \tau(X,L,L;\epsilon) &= \frac{1}{2}\left\lvert  \frac{p^*_H}{p_{L}^*} \right\rvert \sum_{i=1}^I \eta_i |\theta_i - \theta(X)| \epsilon + \mathcal{O}(\epsilon^2).
\end{align}

\end{proof}
